\section{Arithmetic scales and divisor weights}\label{sec:arithmetic}

Prediction compares a count sampled at the pivots with a count sampled
at the block products.  For a block \(B=T\cup\{i\}\), conditioning on
\(\sigma=t_T\) changes the harmonic density of \(t_i\) to
\(\sigma\1_{\sigma\mid y}/(Z_i y)\) on
\([\sigma X_i,\sigma X_i^2)\), with the same \(W\)-unit restriction.
Averaging the divisibility factor gives \(\nu_B\).  We first show that
returning the endpoints to \([X_i,X_i^2)\) has negligible total mass.

We then construct scales and prime parameters for the correlation
argument.  The prime laws have enough harmonic mass to permit prime
insertion, and their residue uniformity makes products of divisor
weights have mean one on the required linear systems.  A separate
coprimality estimate for independent polynomial values will allow
Section~\ref{sec:prediction} to combine cube tests with different
prime-dependent step sizes.

We use \emph{rough} to mean coprime to \(W=\prod_{p\le w}p\), and put
\[
 |a|_{>w}=\prod_{p>w}p^{v_p(a)}\qquad(a\in\Z\setminus\{0\}).
\]
All finite templates in this Section are fixed before the limit
\(w\to\infty\).  A template specifies numbers of variables and rows,
block sizes, and integer polynomials, but no color functions.

\subsection{Harmonic sampling identities}

We give explicit total-mass bounds to distinguish multiplication of
an argument from conditioning on divisibility.  For a signed measure
on \(\Z\), write \(\|\xi\|_1=\sum_{n\in\Z}|\xi(n)|\).
Thus probability total variation is one half of this norm.

\begin{lemma}[Sampling and changes of variables]\label{lem:sampling}
Let \(X\ge2\) be an integer, set
\[
 \vartheta_W=\frac{\phi(W)}W,\quad L_X=\log X,\quad
 Z_X=\sum_{\substack{X\le n<X^2\\(n,W)=1}}\frac1n,\quad
 \mu_X(n)=\frac{\1_{[X,X^2)}(n)\1_{(n,W)=1}}{nZ_X},
\]
and suppose \(L_X>W/X\).  Let \(Y\) have law \(\mu_X\).
\begin{enumerate}
\item For \((k,W)=1\), every residue \(a\pmod k\), and
\(0<A<B\),
\begin{equation}\label{eq:periodic-harmonic}
 \left|\sum_{\substack{A\le n<B\\(n,W)=1\\n\equiv a\ (k)}}\frac1n
  -\frac{\vartheta_W}{k}\log\frac BA\right|
 \le\frac{\phi(W)}A.
\end{equation}
In particular \(|Z_X-\vartheta_WL_X|\le\phi(W)/X\), and
\begin{equation}\label{eq:harmonic-residue-error}
 \left|k\PP(Y\equiv a\pmod k)-1\right|
 \le E_X(k):=\frac{W(k+1)}{X(L_X-W/X)}.
\end{equation}
The total-mass distance of the residue law from uniform measure is
also at most \(E_X(k)\).
\item If \(h\in W\Z\), then
\begin{equation}\label{eq:harmonic-translation}
 \|\operatorname{Law}(Y+h)-\mu_X\|_1
 \le \min\left(2,\frac{2|h|}{X(L_X-W/X)}\right).
\end{equation}
\item If \(1\le k\le X\) and \((k,W)=1\), let
\(\eta_k(n)=k\1_{k\mid n}\mu_X(n)\), an unnormalized positive
measure.  Then
\begin{equation}\label{eq:harmonic-dilation}
 \|\operatorname{Law}(kY)-\eta_k\|_1
 \le\frac{2\log k+(Wk/X)(1+1/X)}{L_X-W/X},
\end{equation}
and \(|\eta_k(\Z)-1|\le E_X(k)\).
\end{enumerate}
For the asymptotic conclusions, let \(K,H,V\ge1\) be auxiliary
quantities depending on \(w\).  Each residue and translation error is
\(o(V^{-C})\), for every fixed \(C\), uniformly for
\(k\le K\) and \(|h|\le H\), if \(X\) dominates every fixed
power of \(2+W+K+H+V\).  The dilation error has the same conclusion
if \(\log X\) dominates every fixed power of \(2+W+K+V\).
These conclusions still hold in expectations of functions bounded
by any fixed power of \(V\).
\end{lemma}

\begin{proof}
The indicator in \eqref{eq:periodic-harmonic} is periodic modulo
\(Wk\) and occupies exactly \(\phi(W)\) residue classes.  On any
interval, its counting function differs from
\(\vartheta_W/k\) times length by at most \(\phi(W)\): apply
the discrepancy bound of one to each occupied class.  Partial
summation against \(1/t\) gives \eqref{eq:periodic-harmonic}.
Apply it with \(k=1\) to obtain the bound on \(Z_X\).
For a general class the error in its unnormalized mass is at most
\(\phi(W)/X\); subtracting \(Z_X/k\) and using
\(Z_X\ge\vartheta_W(L_X-W/X)\) proves
\eqref{eq:harmonic-residue-error}.  Sum the absolute class errors
to obtain the total-mass assertion.

For translation put \(H=|h|\).  Translation invariance of the norm
reduces to \(h=H\ge0\).  On the overlap of the supports,
\(1/(n-H)\ge1/n\), and \(W\mid H\) preserves the unit
condition.  The negative part of
\(\operatorname{Law}(Y+H)-\mu_X\) is therefore precisely the
original mass on \([X,\min(X+H,X^2))\).  Both measures have
mass one, so the norm is twice that mass.  If
\(H\le X^2-X\), the interval \([X,X+H)\) has exactly
\(\vartheta_WH\) units modulo \(W\), since its integer length
is a multiple of \(W\); its harmonic mass is at most
\(\vartheta_WH/X\).  If \(H>X^2-X\), the norm is at most two,
and the second bound in \eqref{eq:harmonic-translation} is already
at least two.  This proves the translation estimate.

For dilation, at a multiple \(n\) of \(k\) the pushed-forward
mass is \(k/(nZ_X)\), supported on \([kX,kX^2)\) and on
\((n,W)=1\).  Thus it agrees exactly with \(\eta_k\) on the
overlap.  Changing variables \(n=ku\) in the two boundary pieces,
and using \(k\le X\), gives the exact formula
\begin{align*}
 \|\operatorname{Law}(kY)-\eta_k\|_1
 =\frac1{Z_X}\bigg(
  \sum_{\substack{X/k\le u<X\\(u,W)=1}}\frac1u
  +\sum_{\substack{X^2/k\le u<X^2\\(u,W)=1}}\frac1u
 \bigg).
\end{align*}
Each main term in \eqref{eq:periodic-harmonic} is
\(\vartheta_W\log k\), and the two errors are at most
\(\phi(W)k/X\) and \(\phi(W)k/X^2\), respectively.
Dividing by the lower bound for \(Z_X\) proves
\eqref{eq:harmonic-dilation}.  Its mass assertion follows from
\eqref{eq:harmonic-residue-error} with \(a=0\).
Finally, integration against a function of supremum norm \(B\)
costs at most \(B\) times the total-mass error.  The stated growth
conditions absorb every fixed choice of \(B=V^{O(1)}\).
\end{proof}

For later use, uniform sampling on an integer interval of length
\(T\) has residue law modulo \(k\) at total-mass distance at most
\(2k/T\) from uniform measure.  Each residue occurs either
\(\lfloor T/k\rfloor\) or \(\lceil T/k\rceil\) times, up to
the harmless endpoint convention.  The same bound is uniform in
the interval's location.  Translating the interval by an integer
\(u\) changes its probability law in total-mass norm by at most
\(2\min(1,|u|/T)\).  Product laws incur the sum of the coordinate
errors, by telescoping their tensor products.  These facts and
Lemma~\ref{lem:sampling} apply conditionally when the interval
endpoints, dilations, or translations have been fixed by outside
variables and the displayed bounds hold uniformly in those variables.

\subsection{Recovering the product sampling law}

The next Corollary identifies the center weights in
\eqref{eq:weighted-count} with the actual block-product sampling.
Once the shifted factors have bounded models, this identity will remove
the remaining weights from the count.

\begin{corollary}[Product law]\label{cor:product-law}
Let \(B=T\cup\{i\}\) be a block, and let \(t_j\) be the
independent harmonic samples at its raw cutoffs.  Let \(V=V(w)\ge1\)
be an auxiliary scale, and put
\(K_T=\prod_{j\in T}X_j^2\).  If \(\log X_i\) dominates
every fixed power of \(2+W+K_T+V\), then
\begin{equation}\label{eq:block-product-law}
 \left\|\operatorname{Law}(t_B)-
                  \nu_B\mu_i\right\|_1=o(V^{-C})
 \qquad\text{for every fixed }C>0.
\end{equation}
For any fixed family of disjoint blocks the corresponding joint
law differs in the same sense from the product of their weighted
pivot measures, when these hypotheses hold at each pivot.
\end{corollary}

\begin{proof}
Condition on \(\sigma=t_T\).  It is coprime to \(W\),
independent of \(t_i\), and bounded by \(K_T\).  Apply
\eqref{eq:harmonic-dilation} with \(k=\sigma\), uniformly in
this range, and average.  The comparison measure is exactly
\[
 \mu_i(y)\E_{\sigma}\sigma\1_{\sigma\mid y}
     =\mu_i(y)\nu_B(y),
\]
proving \eqref{eq:block-product-law}.  In particular its mass is
\(1+o(V^{-C})\).  Disjoint blocks use disjoint raw variables, so
their original laws are independent.  Telescoping a fixed tensor
product bounds its total-mass error by the sum of the individual
errors times the masses of the other factors, all of which are
\(1+o(1)\).  This proves the joint assertion.
\end{proof}

\subsection{A sequential construction of the master scales}

We now work on a master index set \([N]\). Here \(R_l\) are auxiliary gap
lengths and \(X_j\) are cutoffs for the raw variables.
Section~\ref{sec:prediction} will select the final \(n\) indices and choose
each \(H_i\) from these gap lengths.
The smooth factors \(h_j\) make all required shifts integral.  The prime
pools supply the harmonic mass and residue uniformity described above;
the gap lengths absorb the finitely many polynomial divisibilities
needed by the subsequent changes of variables.  The raw cutoffs are
chosen last in each gap so that the preceding sampling estimates apply.

\begin{lemma}[Master scales]\label{lem:master-scales}
Fix \(N\), a bound on block sizes, a finite set
\(\mathcal A\subset\Q_{>0}\), and a finite list \(\mathcal D\) of
nonzero integer polynomials in finitely many prime-parameter slots.
There are choices of \(h_j,M,X_j,R_l,P_l^-,P_l^+\), for
\(1\le j,l\le N\), with the following properties.
\begin{enumerate}
\item
\(h_j=W^{w2^{N-j}}\), and \(M\) is a power of \(W\), divisible by
\(W^w\), with \(h_j\le M\) for every index and \(h_B\le M\)
for every block under consideration.
Every permitted adding pair satisfies
\(h_A/h_B\in W^w\N\).  For a master chain and
\(c_d=h_{B_d}a_d\), \(a_d\in\mathcal A\), all \(c_d\) are integers
and \(c_u/c_d\in W\N\) when \(u<d\), for sufficiently large \(w\).
\item
There is an integer \(e_0=e_0(w)\) such that, for independent uniform
units in the slots modulo \(W^{e_0}\),
\begin{equation}\label{eq:small-prime-exception}
 \PP\bigl(p^{e_0}\mid D(\boldsymbol u)
       \text{ for some }p\le w,\ D\in\mathcal D\bigr)=o(1).
\end{equation}
Moreover \(W^{e_0+1}c_d\mid M\) for every possible \(c_d\).
\item
Writing
\[
 V_l=2+M+\prod_{j<l}X_j^2,
 \qquad
 Q_l=W^{e_0}\prod_{w<p\le V_l}p,
\]
the pool \(\mathcal P_l\) consists of all primes in a finite union of
consecutive complete dyadic intervals
\([P_l^-,P_l^+)\).  Its law and harmonic mass are
\begin{equation}\label{eq:prime-pool-law}
 H_l^{\rm pr}=\sum_{p\in\mathcal P_l}\frac1p,
 \qquad \lambda_l(p)=\frac1{pH_l^{\rm pr}}.
\end{equation}
Both \(P_l^-\) and \(H_l^{\rm pr}\) dominate every fixed power of
\(V_l\).  The residue law of a sample modulo \(Q_l\) differs from
uniform measure on \((\Z/Q_l\Z)^\times\) by
\(o(V_l^{-C})\) in total variation, for every fixed \(C\).
All slots are sampled independently, including slots in different gaps.
\item
\(R_l\) dominates every fixed power of \(P_l^++V_l\).  It is divisible
by \(M\), by every earlier \(R_j\), and by
\(M|D(\boldsymbol p)|\) for all the nonzero values required for
divisibility in gap \(l\), with these polynomial slots drawn from
that gap's pool.  The integer \(X_l\) is a power of two, and \(\log X_l\)
dominates every fixed power of \(R_l\).
\end{enumerate}
One can include any prescribed finite family of polynomial nonvanishing
conditions in \(\mathcal D\).  Repeated entries and zero values then
have probability \(o(V_l^{-C})\) for every fixed \(C\), whereas the
small-prime exceptions in \eqref{eq:small-prime-exception} have
probability \(o(1)\).
\end{lemma}

\begin{proof}
Set \(e_j=2^{N-j}\).  The strict inequality
\(e_j>\sum_{k>j}e_k\) shows that the sign of
\(\sum_{j\in A}e_j-\sum_{j\in B}e_j\), for distinct subsets, is
determined by the smallest index in their symmetric difference.
For an adding pair that index belongs to \(A\); for two chained
blocks it belongs to the earlier block.  The positive difference is
an integer, so the corresponding ratio of \(h\)-products is divisible
by \(W^w\).  Fixed rational multipliers have bounded valuations and
only finitely many denominator primes.  Consequently they are absorbed
by these powers for all sufficiently large \(w\), leaving at least
one factor of every prime dividing \(W\) in each required ratio of
the \(c_d\)'s.

For completeness, a nonzero polynomial over \(\Q_p\) has a null zero
set in a product of \(p\)-adic unit groups.  In one variable it has
only finitely many roots.  In several variables, expand in a variable
on which it depends and choose a nonzero coefficient polynomial.
By induction, the set on which that coefficient vanishes has measure
zero; off that set, each fiber has finitely many roots.  Fubini's
Theorem proves the assertion.  The decreasing events
\(p^e\mid D\) therefore have probabilities tending to zero as
\(e\to\infty\).  At any fixed \(w\), the list of polynomials and
primes is finite.  Choose a common \(e_0\) making their union
probability at most \(1/w\).  Reduction of uniform units modulo
\(W^{e_0}\) has exactly the required distribution modulo
\(p^{e_0}\).  We can now choose the power \(M\) large enough for
all the finitely many stated smooth divisibilities and bounds.

Proceed in increasing order of \(l\).  At this stage \(Q_l\) is a
fixed integer.  The prime number theorem in each reduced arithmetic
progression modulo \(Q_l\) \cite[Equation~(1.1)]{SelbergAP}, followed by partial summation, says that
the harmonic prime law on \([Y,2Y)\) tends to the uniform law on
its unit classes as \(Y\to\infty\).  Since there are finitely many
classes, choose a dyadic lower endpoint so large that the total
variation error is at most \(V_l^{-w}\) on every subsequent dyadic
interval, and also \(P_l^-\ge wV_l^w\).  Increasing the upper
endpoint gives \(H_l^{\rm pr}\ge wV_l^w\), since the harmonic prime
series diverges.  A mixture of measures having the same error bound
has that error bound.  Thus this finite pool has the asserted
residue distribution, however long the union of intervals is.

Only finitely many polynomial values occur in a finite pool.  Take
their nonzero absolute values, the earlier gap lengths, and \(M\),
form the required least common multiple, and choose a multiple
\(R_l\) at least \(w(P_l^++V_l)^w\).  Finally choose a power of two
\(X_l\) with \(\log X_l\ge wR_l^w\).  There is no upper bound in
any of these choices, and no later requirement changes an earlier
choice.

Here is the elementary zero-value estimate used for the final
assertion.  If independent variables each have maximum atom at most
\(a\), a nonzero polynomial of total degree \(d\) vanishes with
probability at most \(da\).  Expand in its last relevant variable:
outside the zero set of its leading coefficient there are at most
as many roots as its degree in that variable, and induction bounds
the leading-coefficient exception by its degree times \(a\).
For our pool, \(a\le1/(P_l^-H_l^{\rm pr})\); the same estimate for
\(x_i-x_j\) handles repeats.  These bounds are smaller than every
fixed inverse power of \(V_l\). Finally, the residue-law approximation
modulo \(Q_l\), together with the Chinese remainder theorem (CRT),
transfers \eqref{eq:small-prime-exception} to the actual prime slots.
\end{proof}

For a block \(B=T\cup\{i\}\) at these scales, its tail bound
\(K_T\) and \(W\) are at most \(V_i\).  Since \(\log X_i\)
dominates every fixed power of \(V_i\),
Corollary~\ref{cor:product-law} applies with \(V=V_i\).
In particular, the joint law comparison holds for every chain and
hence for bounded expectations of all its block variables.

\begin{remark}[Countably many fixed tests]\label{rem:arithmetic-diagonal}
Later, a fixed test may be raised to any fixed moment or combined
with any fixed number of other tests.  This does not require a
moment order growing inside an estimate.  Enumerate the resulting
integer-polynomial templates and positive integer auxiliary parameters.
At stage \(j\) impose the requirements for the first \(j\) templates.
Each stage has finitely many requirements, so the preceding proof
applies.  Choose increasing thresholds \(w_j\) so that the finitely
many error bounds at that stage, including constants in the estimates
below, are at most \(1/j\) once \(w\ge w_j\).  At a given \(w\)
use stage \(\max\{j:w_j\le w\}\), with the growth exponents also
tending to infinity.  Every fixed finite test family then satisfies
all the conclusions eventually.  This diagonal choice can be made
before the functions to be tested are chosen, because its templates
depend only on their algebraic forms.  All moment limits below mean
first \(w\to\infty\) at a fixed moment order.
\end{remark}

\subsection{Coprimality beyond the CRT cutoff}

The next estimate treats prime factors larger than \(V_l\), as well
as those controlled by the residue distribution of a pool.  The
linear-forms estimate below uses residue uniformity separately.
Coprimality has a different purpose: in a cyclic group, the subgroups
of multiples of two coprime integers together generate the whole group.
Section~\ref{sec:prediction} applies this fact to combine cube tests
with different step sizes.

\begin{lemma}[Rough coprimality]\label{lem:rough-coprimality}
Let \(F\) and \(G\) be fixed nonzero integer polynomials in disjoint
tuples of independent prime variables.  Each variable is sampled with
probability proportional to \(1/p\) in a dyadic interval \([Y,2Y)\).
Let \(L\) be the smallest of these lower endpoints.  For all
sufficiently large \(w,L\),
\begin{align}
 \PP(FG=0)&\ll_{F,G}\frac{\log L}{L},\label{eq:polynomial-zero-bound}\\
 \PP\bigl(FG\ne0,\ \gcd(|F|_{>w},|G|_{>w})>1\bigr)
 &\ll_{F,G}\frac1w+\frac{\log L}{\sqrt L}.
 \label{eq:rough-coprimality-bound}
\end{align}
The constants are independent of the ratios between the dyadic
endpoints.  Consequently two independent tuples from any of the
constructed pools have coprime rough polynomial values with
probability \(1-o(1)\).  The same assertion holds after restricting
each tuple to a prime-only event of probability \(1-o(1)\).
\end{lemma}

\begin{proof}
The prime number theorem gives a maximum atom
\(O(\log Y/Y)\) in the harmonic prime law on \([Y,2Y)\).
The polynomial zero estimate in the proof of
Lemma~\ref{lem:master-scales} proves
\eqref{eq:polynomial-zero-bound}, uniformly in the endpoints.

We shall also use the following consequence of the interval
Brun--Titchmarsh inequality.  If \(p\le\sqrt Y\) is prime, a
harmonically sampled prime in \([Y,2Y)\) belongs to any prescribed
nonzero class modulo \(p\) with probability \(O(1/p)\).
Indeed Theorem~2 of \cite{Yamada}, together with
\(\log(Y/p)\ge\tfrac12\log Y\), bounds the number of primes in a
reduced class in that interval by \(O(Y/(p\log Y))\).
Divide by the prime count \(\gg Y/\log Y\).
Harmonic and uniform sampling on that interval
have densities within a factor of two.  The zero class contains no
prime in \([Y,2Y)\) when \(p\le\sqrt Y\).

Remove variables on which the polynomials do not depend, and argue
by induction on the total number of remaining variables.  A fixed
nonzero constant has no prime factors larger than \(w\) once \(w\)
is large, so it gives the base case.  In each nonconstant polynomial
choose as its main variable one whose interval has largest lower
endpoint in that tuple.  Denote those endpoints by \(Y_F,Y_G\),
and the corresponding nonzero leading-coefficient polynomials by
\(A,B\).  The events \(A=0\) or \(B=0\) cost
\(O_{F,G}(\log L/L)\).  If a common prime divisor of \(F,G\)
also divides \(A\), apply the induction hypothesis to \(A,G\);
if it divides \(B\), apply it to \(F,B\).  The two tuples in each
application remain disjoint, and their total number of variables
strictly decreases.  It remains to consider common primes dividing
neither leading coefficient.

Put \(Y_0=\min(Y_F,Y_G)\).  For \(w<p\le\sqrt{Y_0}\), condition
on all variables except the two main variables.  When the respective
leading coefficient is nonzero modulo \(p\), each polynomial has at
most its fixed degree many roots modulo \(p\).  The preceding
Brun--Titchmarsh consequence bounds the two independent tests by
\(O_{F,G}(p^{-2})\).  Their sum is \(O_{F,G}(1/w)\).

For the remaining primes, expose the entire tuple whose main endpoint
is \(Y_0\), say the \(F\)-tuple.  Every variable in that tuple is
at most \(2Y_0\), so, when \(F\ne0\),
\(1\le |F|\le C_FY_0^{\deg F}\).  It follows that \(F\) has
only \(O_F(1)\) distinct prime divisors exceeding \(\sqrt{Y_0}\).
Test each such prime \(p\) against the independent \(G\)-tuple,
whose main endpoint \(Y=Y_G\) satisfies \(Y\ge Y_0\).  If
\(p\le\sqrt Y\), the cost is \(O_G(1/p)=O_G(Y_0^{-1/2})\),
again excluding the already treated event \(p\mid B\).  If
\(p>\sqrt Y\), each of the boundedly many root classes contains
at most \(Y/p+1\) integers of the interval.  Dividing by
\(\gg Y/\log Y\) gives
\[
 O_G\left(\frac{\log Y}{\sqrt Y}\right)
 =O_G\left(\frac{\log Y_0}{\sqrt{Y_0}}\right).
\]
The last inequality holds for all sufficiently large \(Y_0\), since
\(\log Y/\sqrt Y\) is then decreasing.  Summing over the bounded
number of large prime factors proves the induction and
\eqref{eq:rough-coprimality-bound}.

Conditioning pool samples on their individual dyadic intervals leaves
independent harmonic prime laws.  The established bounds are uniform
over these choices, so averaging proves the pool assertion.
Restricting to events of probabilities \(1-o(1)\) changes any
probability by \(o(1)\).
\end{proof}

The prime pools now have both residue uniformity and rough coprimality.
We turn to the mean-one estimate that permits repeated
Cauchy--Schwarz with divisor weights.

\subsection{The weighted linear-forms estimate}

Recall that for a block \(B=T\cup\{i\}\) and independent tail
product \(\sigma=t_T\), the divisor weight is
\begin{equation}\label{eq:arithmetic-divisor-weight}
 \nu_B(y)=\E_{\sigma}\sigma\1_{\sigma\mid y}.
\end{equation}
In expanding a product of weights, every occurrence of a weight
receives a fresh independent draw, even when its block label repeats.
The next Proposition records precisely the row hypotheses needed
below.  In particular, restrictions on prime parameters are permitted.

The mean-one linear-forms estimate below plays the role of the
pseudorandomness conditions in Green and Tao's transference method
\cite[Definition~3.1 and Section~5]{GreenTaoPrimes}.
We prove the estimate for the present divisor weights and
prime-dependent forms.

\begin{proposition}[Weighted linear forms]\label{prop:linear-forms}
Fix the number \(q\) of rows, the number \(d\) of base variables,
and a bound \(b\) on the number of raw factors in each divisor.
For each \(u\in[q]\), let \(\sigma_u\) be a product of at most
\(b\) independent harmonic \(W\)-unit variables at the master
cutoffs of Lemma~\ref{lem:master-scales}, and write
\(\nu_u(y)=\E\sigma_u\1_{\sigma_u\mid y}\).
All these raw draws are independent across occurrences.  Suppose
\(\sigma_u\le V\) surely, where \(V\ge M\to\infty\).
Let \(\boldsymbol p\) consist of independent prime slots from
finitely many gaps, each having \(V_l\ge V\).  These slots are
independent of all divisor draws.

Let \(\mathcal G\) be a domain depending only on
\(\boldsymbol p\), and let
\(\ell_u(\boldsymbol x;\boldsymbol p)\) be homogeneous linear
rows.  On \(\mathcal G\) their arguments are integers on the
base sampling support.  Assume the following conditions.
\begin{enumerate}
\item Given \(\boldsymbol p\) and any possible divisor draws,
the joint residue law of \(\boldsymbol x\) modulo
\(K=\prod_{u=1}^q\sigma_u\) differs in total variation by at
most \(\epsilon_{\rm base}\) from the uniform law on
\((\Z/K\Z)^d\).  The bound is uniform on \(\mathcal G\).
\item For every \(w<p\le V\), each row is nonzero modulo
\(p\) on \(\mathcal G\).  There is a fixed finite list of
nonzero integer polynomials in the prime slots such that, if all
their values are nonzero modulo \(p\), every pair of rows is
linearly independent over \(\mathbb F_p\).
\end{enumerate}
The CRT law of all prime slots at the primes \(w<p\le V\)
is within total variation \(\epsilon_{\rm CRT}\) of independent
uniform units, with independence also across these primes.
Then, uniformly in every further prime-only event
\(\mathcal E\subseteq\mathcal G\),
\begin{equation}\label{eq:restricted-linear-forms}
 \E_{\boldsymbol p,\boldsymbol x}
   \1_{\mathcal E}(\boldsymbol p)
   \prod_{u=1}^q\nu_u(\ell_u(\boldsymbol x;\boldsymbol p))
 =\PP(\mathcal E)+
 O\left(\frac1w+V^q(\epsilon_{\rm base}+
                              \epsilon_{\rm CRT})\right).
\end{equation}
The error is absolute, so the statement also applies when
\(\PP(\mathcal E)\) tends to zero.
The implied constant depends only on the fixed row and divisor
templates and polynomial tests.  In particular the error is \(o(1)\)
with the scales of Lemma~\ref{lem:master-scales} and base residue
errors smaller than every fixed inverse power of \(V\).

Rows may have rational coefficients if their denominators are units
modulo every possible divisor: interpret their residues by inverting
those denominators.  The integer-value and row hypotheses must still
hold.  Clearing such denominators leaves the congruence calculation
unchanged.  In particular this permits smooth denominators and prime
pool denominators larger than \(V\).
\end{proposition}

\begin{proof}
We first express the normalized divisibility count as a product of local
kernel counts.  Their excess over one is nonnegative; this lets us bound
it without retaining the prime-only event.  We then dominate the divisor
draws by laws with independent prime valuations and sum the local excesses.

\paragraph{Local kernel counts.}
Expand each weight with its independent divisor draw.  Replacing the
base residue law by exact uniform measure modulo \(K\) costs at
most \(2\epsilon_{\rm base}\prod_u\sigma_u
\le2V^q\epsilon_{\rm base}\).  The CRT theorem now factors the
normalized divisibility count as
\[
 \prod_{w<p\le V}\alpha_p,
 \qquad
 \alpha_p=p^{\sum_u a_u}
  \PP_{\boldsymbol x\bmod p^{A}}
       (p^{a_u}\mid\ell_u(\boldsymbol x)\text{ for every }u),
\]
where \(a_u=v_p(\sigma_u)\) and \(A=\max_u a_u\).
Primes with all \(a_u=0\) contribute one.  Homogeneity makes the
divisibility conditions the kernel of a homomorphism into
\(\prod_u\Z/p^{a_u}\Z\).  Its image has size at most
\(p^{\sum a_u}\), so \(\alpha_p\ge1\).  With a single
positive valuation, the corresponding primitive row is surjective
modulo every power of \(p\), giving \(\alpha_p=1\).
In general one primitive row with valuation \(A\) bounds the
kernel probability by \(p^{-A}\).  If the two rows having the
largest valuations \(A\ge B>0\) are independent modulo \(p\),
one of their two-column minors is a unit.  Fixing the other
coordinates, the two selected coordinates are uniformly and
bijectively mapped to the pair of row values modulo \(p^A\).
Their required divisibilities consequently have probability
\(p^{-A-B}\).  Additional rows can only reduce the probability.

We bound the nonnegative excess uniformly, so that an arbitrary
prime-only domain can subsequently be retained.  Let \(\mathcal T_p\)
be the test that at least one of the prescribed polynomial values
vanishes modulo \(p\).  Set \(\beta_p=0\) when fewer than two
valuations are positive; otherwise set
\[
 \beta_p=
 \begin{cases}
 p^{\sum a_u-A-B}-1,&\mathcal T_p\text{ fails},\\
 p^{\sum a_u-A}-1,&\mathcal T_p\text{ holds}.
 \end{cases}
\]
On \(\mathcal G\), the preceding kernel bounds give
\(0\le\alpha_p-1\le\beta_p\).  For every choice of the
parameters and the true, bounded divisor draws,
\begin{equation}\label{eq:local-excess-domination}
 0\le\1_{\mathcal E}
       \left(\prod_p\alpha_p-1\right)
 \le\prod_{w<p\le V}(1+\beta_p)-1
 \le\prod_u\sigma_u\le V^q.
\end{equation}
In particular, we can first replace the prime residues by independent
CRT-uniform units at cost at most \(2V^q\epsilon_{\rm CRT}\).
This replacement is made while all divisor draws are still bounded.

\paragraph{Independent comparison laws.}
We next compare those draws with laws whose prime valuations are
independent.  For a raw factor at cutoff \(X_j\), put
\(\epsilon_j=1/\log X_j\) and define
\[
 \widetilde\mu_j(n)
  =\frac{\1_{(n,W)=1}n^{-1-\epsilon_j}}{
    \zeta(1+\epsilon_j)\prod_{p\le w}(1-p^{-1-\epsilon_j})}
 \qquad(n\ge1).
\]
Its denominator is \((1+o(1))\vartheta_W\log X_j\).
Indeed \(\epsilon\zeta(1+\epsilon)\to1\), and
\[
 0\le\log\prod_{p\le w}
        \frac{1-p^{-1-\epsilon}}{1-p^{-1}}
 \le\epsilon\sum_{p\le w}\frac{\log p}{p-1}
 \le\epsilon\log W=o(1).
\]
Lemma~\ref{lem:sampling} gives the same asymptotic normalization
for the original harmonic law on \([X_j,X_j^2)\), and
\(n^{\epsilon_j}\le e^2\) on that interval.  Hence
\(\mu_j(n)\le C_0\widetilde\mu_j(n)\) pointwise, with an
absolute constant \(C_0\) for sufficiently large \(w\).
For at most \(bq\) independent raw draws this incurs one overall
constant \(C_0^{bq}\), rather than a separate constant for each
prime.

The Euler product defining \(\widetilde\mu_j\) shows that its
valuations at different primes are independent, with
\[
 \widetilde\PP(v_p(n)=a)
    =(1-p^{-1-\epsilon_j})p^{-a(1+\epsilon_j)}
    \le p^{-a}.
\]
Thus a product of at most \(b\) raw factors has valuation mass
at \(a\) at most \((a+1)^b p^{-a}\).  This bound is uniform
in the individual cutoffs.  Divisors remain independent of each
other and of all prime parameters.

\paragraph{Summing the local excesses.}
For sufficiently large \(w\), none of the fixed nonzero test
polynomials is identically zero modulo any prime \(p>w\).
The elementary polynomial root bound on the product grid
\((\mathbb F_p^\times)^k\) gives
\(\PP(\mathcal T_p)=O(1/p)\).  Under the CRT-uniform law the
tests at different primes are independent.  Combining the valuation
bounds with the definition of \(\beta_p\), the contribution when
\(\mathcal T_p\) fails is at most a constant times
\[
 \sum_{A\ge B\ge1}
  (A+1)^{C_1}(B+1)^{C_1}p^{-A-B}
 \ll p^{-2}.
\]
To see this sum explicitly, choose the rows with the two largest
valuations; all other valuations lie between zero and \(B\),
and summing their polynomial factors is bounded by a fixed power
of \(B+1\).  The factors \(p^{-\sum a_u}\) cancel the
normalizing power in \(\beta_p\), leaving \(p^{-A-B}\).
On \(\mathcal T_p\) the same argument, with all lesser
valuations bounded by \(A\), gives
\[
 \sum_{A\ge1}(A+1)^{C_2}p^{-A}\ll p^{-1}.
\]
Its additional \(O(1/p)\) test probability again gives
\(\E\beta_p\ll p^{-2}\).  The two displayed series bounds
follow by factoring out respectively \(p^{-2}\) and \(p^{-1}\)
and bounding the remaining geometrically convergent series using
\(p\ge2\).

We have now obtained independence across primes for every quantity
in the nonnegative right side of
\eqref{eq:local-excess-domination}.  Consequently its expectation,
after the single global divisor-domination constant, is at most
\[
 C_0^{bq}\left\{\prod_{w<p\le V}(1+C_3p^{-2})-1\right\}
 \ll\sum_{p>w}p^{-2}\ll\frac1w.
\]
The original local product has lower bound one on
\(\mathcal E\).  Integrating that lower bound and the estimated
excess, and restoring the two residue-approximation errors, proves
\eqref{eq:restricted-linear-forms}.  Rational denominators which
are units modulo the divisors induce the same homomorphisms after
inversion, so the proof covers the stated rational-row extension.
\end{proof}

Several consequences of Proposition~\ref{prop:linear-forms} will be
used without altering its hypotheses.  A fixed product of factors
\(1+\nu_u\) has main term \(2^q\PP(\mathcal E)\), by expansion
over subsets of rows.  If at least one factor is \(\nu_u-1\),
and the remaining factors are each \(\nu_v\), \(1+\nu_v\),
or \(\nu_v-1\), its average is \(o(1)\), provided every
subsystem in that expansion meets the same row conditions.  This
is cancellation of the constant main terms.  In particular a
prime-only exceptional event of probability \(o(1)\), contained
in a domain of primitive rows, has weighted mass \(o(1)\).
One may also condition on a prime-only event whose probability is
bounded below, dividing \eqref{eq:restricted-linear-forms} by that
probability.

These statements include systems assembled from different gaps.
Use the common divisor bound \(V\); every relevant pool has
\(V_l\ge V\), so projecting its CRT law supplies the required
accuracy.  Independent replicas give independent slots.  For
coefficients involving
\(M(\boldsymbol p)=M|D(\boldsymbol p)|_{>w}\), a prime
\(w<p\le V\) dividing such a coefficient must divide
\(D(\boldsymbol p)\), so the indicated polynomial tests cover
this possible loss of a minor.  Fixed integer contents are
\(w\)-smooth eventually.  Primitivity and separation of each
actual row system still have to be checked; those checks are given
in Sections~\ref{sec:correlation} and \ref{sec:prediction}.

To verify the base-law hypothesis in those applications, every
divisor product \(K\) is at most \(V^q\).  Uniform auxiliary
intervals of lengths dominating all powers of \(V\), and harmonic
pivot variables whose cutoffs dominate all powers of \(W+V\),
are therefore jointly uniform modulo \(K\), with error smaller
than every fixed inverse power of \(V\).  Their conditional
interval locations do not matter.  A short translation changes
such a box by the sum of the displacement-to-side-length ratios;
the translation and dilation bounds above cover harmonic variables.
For example,
\(R_l/(J_0M(\boldsymbol p))\) dominates all powers of the common
tail bound for each fixed \(J_0\), since
\(M(\boldsymbol p)\) has a fixed-power bound in \(P_l^++V_l\).
Taking square roots of these lengths preserves that domination and
gives negligible relative translations.  These observations supply
uniform bounds for the shifted boxes used later, as well as for
unshifted boxes.
